% ==============================================================================
% SECTION 2 — RELATED WORK
% ==============================================================================
\section{Related Work}
\label{sec:related}

\paragraph{LLMs in game-theoretic settings.}
The intersection of \LLM{}s and game theory has developed rapidly into a
coherent subfield \citep{Sun2025_survey}.
\citet{Aher2023_using} used \LLM{}s to replicate human subject behaviour in
behavioural economics experiments. \citet{Brookins2023_playing} and
\citet{Guo2023_gpt} investigated whether \LLM{}s approximate Nash-rational play,
finding mixed evidence across game types. \citet{Akata2025_repeated}
studied repeated normal-form games, while \citet{Fan2024_rational} demonstrated
systematic failures when \LLM{}s are prompted to output individual game actions.
These limitations motivate our strategy-generation approach, which operates at a
higher level of abstraction and separates policy specification from action
execution.

\paragraph{LLMs in social dilemmas.}
\citet{Yocum2023_mitigating} and \citet{Piatti2024_cooperate}
studied \LLM agent behaviour in Markov social dilemmas. \citet{Park2023_generative}
introduced generative agents for simulating social behaviour.
\citet{DeZarza2023_emergent} explored emergent cooperation in
\LLM-extended coevolutionary theory. \citet{Leibo2017_multiagent} grounded such
studies in multi-agent reinforcement learning. Our contribution to this strand is
the analysis of \emph{evolutionary population dynamics}---how strategy types rise
or collapse across generations---rather than individual agent decisions at a
single time step.

\paragraph{The Willis et al.\ (2025) benchmark.}
\citet{Willis2025_llm_ipd} introduced the framework we extend. They prompt
\LLM{}s to generate natural-language \IPD strategies (converted to Python), then
simulate Moran evolutionary processes over populations of attitude-agents
(Aggressive, Cooperative, Neutral). Key findings include a cooperative bias in
both ChatGPT-4o and Claude~3.5~Sonnet, with Claude exhibiting greater difficulty
generating effective aggressive strategies and notable fragility under action
noise. Applying the Self-Refine framework \citep{Madaan2023_self_refine} to strategy
generation narrowed this gap \citep{Willis2025_llm_ipd}, demonstrating that
iterative self-critique can inadvertently enable aggressive strategies, a result
that complicates alignment-by-design assumptions.

\paragraph{Concurrent and subsequent work.}
Three papers have since tested related questions with different frameworks.
\citet{Payne2025_strategic} run evolutionary \IPD tournaments across GPT, Gemini,
and Claude frontier models using variable termination probabilities and 32,000
prose rationales, without a Moran process or attitude-agent populations; they
find the same provider fingerprint we document (Gemini aggressive, OpenAI
cooperative), providing independent corroboration for our cross-provider
divergence finding (H3, Section~\ref{sec:discussion}) via a mechanistically
distinct design. \citet{Vallinder2024_cultural} apply cross-provider evolutionary
selection to the Donor Game under indirect reciprocity with older models
(Claude~3.5~Sonnet, Gemini~1.5~Flash, GPT-4o) and find the ranking
Claude $>$ Gemini $>$ GPT; our GPT-5.4~Mini emerges as the most cooperative of
the four models we test, inverting that OpenAI position and suggesting that model
generation, not just provider, shapes cooperative disposition.
\citet{Willis2026_collective} scale the Willis et al.\ framework to hundreds of
agents under cultural evolutionary dynamics and find that newer \LLM{}s
\emph{worsen} societal outcomes in their setting---an apparent tension with our
result that cooperative bias persists under Moran selection at $n{=}12$. Whether
the selection mechanism (pairwise Moran versus cultural evolution) or population
scale reconciles the divergence is an open question we flag for future work.

\paragraph{Cross-generational and cross-provider comparisons.}
Prior capability comparisons across \LLM providers focused on general reasoning
benchmarks---MMLU \citep{Hendrycks2021_mmlu}, HumanEval \citep{Chen2021_humaneval},
and similar---where cooperative dispositions are absent by design. The work above establishes that provider-level differences emerge in
evolutionary social dilemmas, but each study either uses older model generations
or a single dilemma type. Within-lineage longitudinal analysis---whether
successive generations from the same provider preserve or reverse established
cooperation patterns---remains absent from the literature. Our study targets that
gap by pairing current frontier models against their immediate predecessors under
the same Moran \IPD framework.

\paragraph{Evolutionary game theory foundations.}
Our simulations rest on the Moran process \citep{Moran1958_random}, the classical
model of evolutionary selection in finite populations; \citet{Nowak2004_cooperation}
established its connection to cooperation in \IPD games in finite populations, and
\citet{Traulsen2006_fixation} derived fixation probabilities under this process.
The \IPD framework follows \citet{Axelrod1984_evolution,Axelrod1981_evolution};
\citet{Nowak2006_evolutionary} provides the backbone for convergence analysis in
finite evolutionary games. The noise mechanism follows
\citet{WuAxelrod1995_noise} and \citet{WahlNowak1999_noise}.
These threads---\LLM strategic behaviour, multi-agent social dilemmas, and
evolutionary population dynamics---jointly define the space this paper occupies:
a cross-provider benchmark for next-generation \LLM{}s, including within-lineage
longitudinal comparisons, evaluated under Darwinian selection over \IPD strategy
populations.
